% GENERATED FILE. Edit canonical lessons, then run npm run generate:books.
% canonical-source-hash: fnv1a64:9680f6e2fc7bec6d
% canonical-lessons: TE-C189-poti, TE-C189-vijeta, TE-C189-patakam, TE-C189-abhiruci, TE-C189-navala

\chapter{Competition, Winner, Medal, Hobby, Novel}
\label{ch:te-competition-winner-medal-hobby-novel}
\hlchaptermodality{189}

\begin{chapteropening}
\textbf{By the end of this chapter:} \emph{I can say a competition, a winner, a medal, a hobby and a novel.}

Complete the last lesson of chapter 189: I can say a competition, a winner, a medal, a hobby and a novel.
\end{chapteropening}

\section[\texorpdfstring{\te{పోటీ} (pōṭī)}{పోటీ (pōṭī)}]{\te{పోటీ} (pōṭī) — a competition, a match}

\label{lesson:TE-C189-poti}



\begin{quote}
Before the new one: say the Telugu for a bindi, then the Telugu for henna.
\end{quote}

\subsection*{You'll want to know: \te{పోటీ}}
\textbf{\te{పోటీ}} — \emph{pōṭī} — \textquotedblleft{}a competition, a match\textquotedblright{}.

\begin{grammarlens}[title={What you've built}]
One more word for this chapter.
\end{grammarlens}

\subsection*{Your turn}
\begin{itemize}
\raggedright
  \item \emph{Say it:} \emph{pōṭī}
  \item \emph{Say it:} \emph{pōṭī}, once more
  \item \emph{Say it:} say \emph{pōṭī}
  \item \emph{Recall:} say the Telugu for a bindi, then the Telugu for henna, then say \emph{pōṭī} again
\end{itemize}

\begin{tcolorbox}[breakable,colback=teal!4,colframe=teal!35!black,title={Before you move on}]
What does \te{పోటీ} mean? (\textquotedblleft{}A competition, a match\textquotedblright{}.) Say it once more.
\end{tcolorbox}
\section[\texorpdfstring{\te{విజేత} (vijēta)}{విజేత (vijēta)}]{\te{విజేత} (vijēta) — a winner}

\label{lesson:TE-C189-vijeta}



\begin{quote}
Before the new one: say the Telugu for henna, then the Telugu for a competition.
\end{quote}

\subsection*{You'll want to know: \te{విజేత}}
\textbf{\te{విజేత}} — \emph{vijēta} — \textquotedblleft{}a winner\textquotedblright{}.

\begin{grammarlens}[title={What you've built}]
One more word for this chapter.
\end{grammarlens}

\subsection*{Your turn}
\begin{itemize}
\raggedright
  \item \emph{Say it:} \emph{vijēta}
  \item \emph{Say it:} \emph{vijēta}, once more
  \item \emph{Say it:} say \emph{vijēta}
  \item \emph{Recall:} say the Telugu for henna, then the Telugu for a competition, then say \emph{vijēta} again
\end{itemize}

\begin{tcolorbox}[breakable,colback=teal!4,colframe=teal!35!black,title={Before you move on}]
What does \te{విజేత} mean? (\textquotedblleft{}A winner\textquotedblright{}.) Say it once more.
\end{tcolorbox}
\section[\texorpdfstring{\te{పతకం} (patakaṁ)}{పతకం (patakaṁ)}]{\te{పతకం} (patakaṁ) — a medal}

\label{lesson:TE-C189-patakam}



\begin{quote}
Before the new one: say the Telugu for a competition, then the Telugu for a winner.
\end{quote}

\subsection*{You'll want to know: \te{పతకం}}
\textbf{\te{పతకం}} — \emph{patakaṁ} — \textquotedblleft{}a medal\textquotedblright{}.

\begin{grammarlens}[title={What you've built}]
One more word for this chapter.
\end{grammarlens}

\subsection*{Your turn}
\begin{itemize}
\raggedright
  \item \emph{Say it:} \emph{patakaṁ}
  \item \emph{Say it:} \emph{patakaṁ}, once more
  \item \emph{Say it:} say \emph{patakaṁ}
  \item \emph{Recall:} say the Telugu for a competition, then the Telugu for a winner, then say \emph{patakaṁ} again
\end{itemize}

\begin{tcolorbox}[breakable,colback=teal!4,colframe=teal!35!black,title={Before you move on}]
What does \te{పతకం} mean? (\textquotedblleft{}A medal\textquotedblright{}.) Say it once more.
\end{tcolorbox}
\section[\texorpdfstring{\te{అభిరుచి} (abhiruci)}{అభిరుచి (abhiruci)}]{\te{అభిరుచి} (abhiruci) — a hobby; taste}

\label{lesson:TE-C189-abhiruci}



\begin{quote}
Before the new one: say the Telugu for a winner, then the Telugu for a medal.
\end{quote}

\subsection*{You'll want to know: \te{అభిరుచి}}
\textbf{\te{అభిరుచి}} — \emph{abhiruci} — \textquotedblleft{}a hobby; taste\textquotedblright{}.

\begin{grammarlens}[title={What you've built}]
One more word for this chapter.
\end{grammarlens}

\subsection*{Your turn}
\begin{itemize}
\raggedright
  \item \emph{Say it:} \emph{abhiruci}
  \item \emph{Say it:} \emph{abhiruci}, once more
  \item \emph{Say it:} say \emph{abhiruci}
  \item \emph{Recall:} say the Telugu for a winner, then the Telugu for a medal, then say \emph{abhiruci} again
\end{itemize}

\begin{tcolorbox}[breakable,colback=teal!4,colframe=teal!35!black,title={Before you move on}]
What does \te{అభిరుచి} mean? (\textquotedblleft{}A hobby; taste\textquotedblright{}.) Say it once more.
\end{tcolorbox}
\section[\texorpdfstring{\te{నవల} (navala)}{నవల (navala)}]{\te{నవల} (navala) — a novel}

\label{lesson:TE-C189-navala}



\begin{quote}
Before the new one: say the Telugu for a medal, then the Telugu for a hobby.
\end{quote}

\subsection*{You'll want to know: \te{నవల}}
\textbf{\te{నవల}} — \emph{navala} — \textquotedblleft{}a novel\textquotedblright{}.

\begin{grammarlens}[title={What you've built}]
One more word for this chapter.
\end{grammarlens}

\subsection*{Your turn}
\begin{itemize}
\raggedright
  \item \emph{Say it:} \emph{navala}
  \item \emph{Say it:} \emph{navala}, once more
  \item \emph{Say it:} say \emph{navala}
  \item \emph{Recall:} say the Telugu for a medal, then the Telugu for a hobby, then say \emph{navala} again
\end{itemize}

\begin{tcolorbox}[breakable,colback=teal!4,colframe=teal!35!black,title={Before you move on}]
What does \te{నవల} mean? (\textquotedblleft{}A novel\textquotedblright{}.) Say it once more.
\end{tcolorbox}
